\section{Metodología}
\label{ch:metodologia}

\subsection{Propuesta de solución}
\label{subsec:propuesta_solucion}

\subsubsection{Análisis del Contexto y Problemática}
\label{subsubsec:analisis_problematica}

En los juzgados del Poder Judicial del Estado de Tamaulipas (PJETAM), la transcripción y redacción manual de acuerdos judiciales representa un cuello de botella logístico crítico. Aunque los expedientes se digitalizan, suelen almacenarse como imágenes estáticas sin indexación semántica. Esto obliga al personal a leer visualmente cada promoción, buscar artículos legales manualmente y redactar borradores desde cero. 

Esta excesiva carga cognitiva ralentiza la impartición de justicia, eleva el estrés laboral y aumenta el riesgo de errores por fatiga visual, haciendo que el modelo operativo actual sea insostenible ante el volumen diario de demandas.

\subsubsection*{Formulación de la Solución}

Para resolver esta crisis, se propone el desarrollo del \textbf{Sistema de Apoyo a la Generación de Acuerdos Judiciales mediante Inteligencia Artificial Generativa (SAGAI)}. Esta solución se materializa en una plataforma con una interfaz web dinámica e interactiva, diseñada para transformar el flujo de trabajo tradicional en un proceso asistido y eficiente.

SAGAI funciona como un intermediario inteligente: automatiza las tareas operativas (extracción de texto y mecanografía) para que el servidor público concentre su experiencia exclusivamente en la validación jurídica y edición final del documento desde su navegador.

\subsubsection{Especificación de Requerimientos del Sistema}
\label{subsubsec:requerimientos}

Para garantizar el éxito de SAGAI, se han definido los siguientes requerimientos estructurales:

\textbf{Requerimientos Funcionales:}
\begin{itemize}
    \item El sistema debe permitir la carga de archivos en formatos estándar de digitalización (PDF, JPG, PNG).
    \item El sistema debe extraer automáticamente el texto de los documentos cargados mediante reconocimiento óptico.
    \item El sistema debe comunicarse con un Modelo de Lenguaje Grande para enviar el texto extraído y recibir una propuesta de acuerdo fundamentada.
    \item El sistema debe proveer un editor de texto interactivo para que el usuario modifique, apruebe y formatee el documento final.
\end{itemize}

\textbf{Requerimientos No Funcionales:}
\begin{itemize}
    \item \textbf{Seguridad:} Los documentos procesados no deben almacenarse de forma permanente en servidores externos no autorizados.
    \item \textbf{Disponibilidad:} La plataforma web debe ser accesible desde la red interna del PJETAM.
    \item \textbf{Rendimiento:} El tiempo de procesamiento entre la carga del documento y la generación del borrador debe ser inferior al tiempo promedio de redacción manual.
\end{itemize}

\subsubsection*{Arquitectura Tecnológica y Pila de Desarrollo}

La plataforma SAGAI se fundamenta en una arquitectura desacoplada y escalable, seleccionando tecnologías robustas para cada capa del sistema:

\textbf{Desarrollo Frontend:}\\
Para la interacción con el usuario, el desarrollo Frontend se centra en la construcción de interfaces dinámicas e interactivas utilizando JavaScript. Se prioriza la manipulación eficiente del Modelo de Objetos del Documento (DOM), el consumo asíncrono de APIs y el manejo del estado de la aplicación para el editor de texto enriquecido. Esto permite un renderizado rápido de los borradores y mantiene una separación estricta y limpia de la lógica de presentación en el navegador.

\textbf{Desarrollo Backend y Visión Computacional:}\\
El núcleo del servidor (\textit{Backend}) será programado en Python, utilizando el marco de trabajo \textit{FastAPI} para asegurar un procesamiento concurrente y de alto rendimiento. En esta capa reside el módulo de Visión Computacional, impulsado por \textit{Tesseract OCR}, responsable de ejecutar la extracción de texto local antes de cualquier transmisión de datos.

\textbf{Gestión de Datos y Entorno de Despliegue:}\\
Los metadatos estructurales y el registro de auditoría de los acuerdos generados serán gestionados mediante un sistema de bases de datos relacional utilizando MySQL. Todo el ecosistema será desplegado sobre un entorno de servidor basado en Linux, específicamente sobre una distribución Debian, lo que garantiza un control administrativo total sobre la personalización del sistema, seguridad a nivel de puertos y una alta estabilidad para ejecutar el motor OCR en producción.

\textbf{Capa Cognitiva (Inteligencia Artificial):}\\
El análisis semántico asistido se delegará a la API de Inteligencia Artificial de Anthropic Claude. Esta capa está configurada bajo la arquitectura de Generación Aumentada por Recuperación (RAG) para anclar y restringir sus sugerencias estrictamente a los Códigos Procesales del Estado de Tamaulipas.

\subsubsection*{Evaluación de Alternativas y Viabilidad}

Durante la fase de diseño, se descartaron enfoques tradicionales por sus limitantes técnicas. La automatización mediante plantillas estáticas fue rechazada por su rigidez, ya que requiere reprogramación ante cualquier reforma de ley o formato inusual de promoción. Asimismo, delegar el procesamiento a servicios comerciales 100\% en la nube sin un filtro previo representaba un riesgo para la privacidad de los expedientes.

Por ello, SAGAI emplea un enfoque híbrido: limpieza local complementada con inferencia de lenguaje en la nube con contexto controlado. Desde la perspectiva de viabilidad institucional, el proyecto resulta altamente rentable. La pila tecnológica base (JavaScript, Python, MySQL, Debian) es de código abierto. Por su parte, la infraestructura cognitiva (API de Anthropic Claude) será gestionada y financiada en su totalidad por el PJETAM a través de las gestiones del asesor empresarial, garantizando que el tribunal asuma únicamente el costo por uso real, eliminando la dependencia de licencias de software prohibitivas.

\subsubsection{Plan de Ejecución y Cronograma Operativo}
\label{subsubsec:plan_ejecucion}

La construcción de la plataforma se ejecutará bajo una metodología de desarrollo modular e incremental. Este enfoque aísla la complejidad del sistema; es decir, permite desarrollar y probar de forma independiente el motor de extracción OCR, el generador de IA y el \textit{Frontend}. 

El desarrollo está planificado para completarse en 14 semanas. A continuación, se describen las fases operativas:

\begin{itemize}
    \item \textbf{Fase 1: Requerimientos y Entorno.} Involucra el levantamiento de procesos con el personal, el diseño de la arquitectura y la configuración del servidor Debian y la base de datos MySQL.
    \item \textbf{Fase 2: Motor OCR.} Centrada en la implementación de la visión computacional y la calibración algorítmica de Tesseract para los formatos de los juzgados.
    \item \textbf{Fase 3: Capa Cognitiva.} Se integra la API de Anthropic Claude y se construye el sistema RAG utilizando los códigos procesales.
    \item \textbf{Fase 4: Desarrollo Frontend.} Programación de las interfaces web y consumo de las APIs del Backend para consolidar la interacción del usuario.
    \item \textbf{Fase 5: Integración y Pruebas.} Ejecución de pruebas de carga, validación del modelo \textit{Human-in-the-loop} y ajustes de seguridad.
    \item \textbf{Fase 6: Cierre.} Compilación de manuales técnicos y entrega formal del proyecto.
\end{itemize}

Para la visualización temporal de estas fases, se presenta el siguiente cronograma condensado:

\begin{table}[H]
    \centering
    \caption{Cronograma condensado de implementación de SAGAI.}
    \label{tab:cronograma_sagai}
    \begin{tabular}{|l|l|l|}
        \hline
        \textbf{Periodo} & \textbf{Semanas} & \textbf{Fase de Desarrollo} \\
        \hline
        Mayo & Semanas 1 - 2 & Fase 1: Requerimientos y Entorno Backend \\
        \hline
        Mayo - Junio & Semanas 3 - 5 & Fase 2: Motor de Extracción Local (OCR) \\
        \hline
        Junio & Semanas 6 - 8 & Fase 3: Capa Cognitiva e Integración IA (Claude) \\
        \hline
        Junio - Julio & Semanas 9 - 10 & Fase 4: Desarrollo Frontend e Interfaces \\
        \hline
        Julio & Semanas 11 - 13 & Fase 5: Pruebas de Integración y Validación \\
        \hline
        Agosto & Semana 14 & Fase 6: Documentación Técnica y Cierre \\
        \hline
    \end{tabular}
    \caption*{\footnotesize \textbf{Fuente: Elaboración propia.}}
\end{table}